\documentclass{article}
\usepackage{amsmath,amssymb,xurl}
\usepackage[hidelinks]{hyperref}
\setlength{\emergencystretch}{2em}
% Shared commands for notes in docs/paper-gaps/.

\DeclareUnicodeCharacter{2080}{\ifmmode{}_{0}\else\textsubscript{0}\fi}
\DeclareUnicodeCharacter{2081}{\ifmmode{}_{1}\else\textsubscript{1}\fi}
\DeclareUnicodeCharacter{2082}{\ifmmode{}_{2}\else\textsubscript{2}\fi}
\DeclareUnicodeCharacter{2099}{\ifmmode{}_{n}\else\textsubscript{n}\fi}

\newcommand{\Lean}{\textsc{Lean}}
\ExplSyntaxOn
\NewDocumentCommand{\leanid}{m}{
  \ifmmode
    \textsf{\detokenize{#1}}
  \else
    \group_begin:
    \urlstyle{sf}
    \tl_set:Nn \l_tmpa_tl {#1}
    \tl_replace_all:Nnn \l_tmpa_tl {\_} {_}
    \exp_args:NV \path \l_tmpa_tl
    \group_end:
  \fi
}
\ExplSyntaxOff
\providecommand{\tr}{\operatorname{tr}}
\newcommand{\tnleanIssueBase}{https://github.com/LionSR/TNLean/issues/}
\newcommand{\tnleanPrBase}{https://github.com/LionSR/TNLean/pull/}
\newcommand{\tnleanDocBase}{https://sirui-lu.com/TNLean/docs/}
\newcommand{\ghissue}[1]{\href{\tnleanIssueBase#1}{GitHub issue \##1}}
\newcommand{\ghpr}[1]{\href{\tnleanPrBase#1}{PR \##1}}
\newcommand{\doclink}[1]{\href{\tnleanDocBase#1.html}{\path{#1}}}

% Dirac notation and the identity operator, as in blueprint/src/macros/common.tex.
% \providecommand keeps note-local definitions authoritative.
\usepackage{dsfont}
\providecommand{\ket}[1]{|#1\rangle}
\providecommand{\bra}[1]{\langle #1|}
\providecommand{\braket}[2]{\langle #1 | #2 \rangle}
\providecommand{\ketbra}[2]{|#1\rangle\!\langle #2|}
\providecommand{\Id}{\mathds{1}}
\providecommand{\one}{\mathds{1}}
\providecommand{\C}{\mathbb{C}}
\providecommand{\MN}[1]{M_{#1}(\C)}
\providecommand{\MD}{M_D(\C)}

% Verdict marker: \gapnote{<kind>}{<status>}. Kind is the policy
% classification (clarification, local-correction, scope-restriction,
% unfaithful, false-source, open-gap); status is open, wip (elimination
% underway), resolved, or historical. Machine-read by the site index;
% prints a verdict line.
\newcommand{\gapnote}[2]{\par\noindent\textbf{Verdict:} #1 (#2).\par}

\title{Normalized error reduction in the final ALVV iteration}
\author{}
\date{November 2, 2026}
\begin{document}
\maketitle
\gapnote{local-correction}{open}

\paragraph{Source.}
The final stage in Arad, Landau, Vazirani and Vidick,
\cite[Figure~fig:algo1, lines 1634--1640]{Arad2016RigorousRG}, alternates
$K=I-H/\lVert H\rVert$ with collective trimming, then selects the $r$
smallest eigenvectors of the compression to the last retained space.
The displayed estimate in the proof of Theorem~t:alg1, line 1813,
asserts, in particular,
\begin{equation}\label{eq:printed-gain}
 \frac{f}{1-\gamma(1-f)/B}
 \ge f\left(1+\frac{\gamma}{2B}\right),
 \qquad f=\alpha^2,\quad B=\lVert H\rVert.
\end{equation}
This inequality fails within the stated range $\alpha\ge0.9$.

\paragraph{An exact nondegenerate example.}
On two qubits take
\[
 H=\operatorname{diag}(0,1,1,1),\qquad
 \psi=\lvert00\rangle,\qquad
 v=\sqrt{99/100}\,\lvert00\rangle+\tfrac1{10}\lvert01\rangle.
\]
The vector $v$ is normalized, $H$ is a positive nearest-neighbour
interaction bounded by the identity, and its ground space is the line
spanned by $\psi$. Its gap and norm are both one. The candidate line
spanned by $v$ is $1/100$-close to this ground space, hence meets the
$3/200$ root-error condition. For requested error $\eta=1/1000$,
\[
 0.9<\sqrt{99/100}<1-\eta,\qquad 99/100<1-\eta.
\]
Thus the example lies before either the amplitude or squared-overlap
accuracy threshold. Since $K=I-H$ sends $v$ to the ground line, the new
normalized squared overlap is one. But the two sides of
\eqref{eq:printed-gain} are, respectively,
\[
 1\qquad\hbox{and}\qquad \frac{99}{100}\frac32=\frac{297}{200}>1.
\]
All vectors and both coefficients in the displayed example are nonzero.
The example concerns the intermediate gain estimate. The conclusion of
Theorem~t:alg1 remains a separate claim requiring a corrected argument.

\paragraph{Normalized error contraction.}
Let a filter preserve a target space and its orthogonal complement,
expand norms on the target, and contract squared norms on the complement
by $\Delta\ge0$.
If a candidate is $\delta$-close to the target with $0\le\delta<1$, its
image is at most
\[
 \delta' = \frac{\Delta\delta}{1-\delta+\Delta\delta}
\]
away from the target. When $\Delta<1$, this contracts the ratio of error
to success: $\delta'/(1-\delta')=\Delta\delta/(1-\delta)$.
The estimate follows by choosing a projection preimage of squared norm
at most $1/(1-\delta)$ and applying Cauchy--Schwarz to its actual filtered
image. It is compatible with fidelities tending to one.

\paragraph{Trimming and extraction.}
Write $\Delta=1-c$, with $0<c\le1$, and fix $0<\eta\le1$.
For $0\le\delta\le1/2$, weighted
comparison with a Schmidt-rank reference and the corrected collective
trimming constant give the sufficient budgets
\[
 \varepsilon_{\rm ref}\le c^2\eta/1024,\qquad
 \xi\sqrt{mb}\le c\eta/48,
\]
and the stable upper recurrence
\[
 \delta_{j+1}\le(1-c/4)\delta_j+c\eta/8.
\]
Here $m$ is the number of actual cuts and $b$ is the rank bound for every
vector in the reference space. These explicit budgets describe one
corrected argument; comparison with every printed quantitative choice
remains required. The normalized estimates alone do not establish the
printed value of the round count.

For the final spectral extraction, closeness of the candidate to the
ground space must also give a trial space of dimension $r$ with the
required energy bound. If $0\le\rho<1$ and $P_W$ projects onto a $\rho$-close candidate,
$G$ is the ground space, and $A=\lVert H-E_0I\rVert$, the actual space
$P_WG$ has dimension $r$ and energy excess at most $A\rho/(1-\rho)$.
This provides the variational comparison for the $r$ smallest compressed
eigenvectors. The actual repeated trims, their parameter choices and
this spectral extraction must be combined before the final source
algorithm can be marked complete.

\bibliographystyle{unsrt}
\bibliography{references}
\end{document}
